% Body file for Section: Perturbative Finiteness (G1)
% Inputted from thqg_perturbative_finiteness.tex
%
% Six stages:
%   (1) HS-sewing framework
%   (2) Shadow free energies and Bernoulli generating function
%   (3) Convergence of the gravitational partition function
%   (4) Heisenberg partition function as prototype
%   (5) Finiteness for the standard landscape
%   (6) Implications for 3d quantum gravity

% ======================================================================
%
%  STAGE 1: THE HS-SEWING FRAMEWORK
%
% ======================================================================

\subsection{The HS-sewing framework}
\label{subsec:thqg-I-hs-sewing-framework}
\index{HS-sewing!perturbative finiteness framework|textbf}
\index{Hilbert--Schmidt!sewing condition}

\begin{definition}[Graded pre-Hilbert core]%
\label{def:g1-graded-pre-hilbert}%
\index{graded pre-Hilbert core|textbf}%
A \emph{graded pre-Hilbert core} for a positive-energy chiral
algebra~$\cA$ is a $\mathbb{Z}_{\geq 0}$-graded vector space
$H = \widehat{\bigoplus}_{n \geq 0} H_n$ equipped with a positive
definite Hermitian form $\langle -, - \rangle_n$ on each weight
space~$H_n$, such that the pair-of-pants multiplication
$\mu \colon \cA \otimes \cA \to \cA$ restricts to bounded
\emph{sector maps}
\begin{equation}\label{eq:g1-sector-maps}
m_{a,b}^c \;\colon\; H_a \,\widehat{\otimes}\, H_b \;\longrightarrow\; H_c,
\qquad
a + b - c \;\in\; \mathbb{Z}_{\geq 0},
\end{equation}
encoding the OPE coefficients $C^{c,k}_{a,i;\,b,j}$ in the
orthonormal bases $\{e_{n,i}\}_{i=1}^{\dim H_n}$.
\end{definition}

\begin{definition}[HS-sewing condition]%
\label{def:g1-hs-sewing}%
\index{HS-sewing!condition|textbf}%
A positive-energy chiral algebra $\cA$ with graded pre-Hilbert
core $(H, \langle -, - \rangle)$ satisfies \emph{HS-sewing}
if there exists $0 < q < 1$ such that the triple series
\begin{equation}\label{eq:g1-hs-sewing}
\Sigma_{\HS}(q)
\;:=\;
\sum_{a,b,c \geq 0} q^{a+b+c}\,\|m_{a,b}^c\|_{\HS}^2
\;<\;\infty
\end{equation}
converges, where $\|m_{a,b}^c\|_{\HS}^2 :=
\sum_{i,j,k} |C^{c,k}_{a,i;\,b,j}|^2$ is the Hilbert--Schmidt
norm of the sector map.  The \emph{weighted completion} is
$H_q := \widehat{\bigoplus}_{n \geq 0} q^n H_n$.
\end{definition}

\begin{proposition}[Closed amplitudes are trace class;
\ClaimStatusProvedElsewhere]%
\label{prop:g1-trace-class}%
\index{trace class!closed amplitudes}%
Under HS-sewing, the weighted multiplication
$m_q \colon H_q \,\widehat{\otimes}\, H_q \to H_q$
is Hilbert--Schmidt.  Every genus-$g$ closed amplitude---obtained
by sewing $2g - 2 + n$ pairs of pants along $3g - 3 + n$
circles---is trace class.
\end{proposition}

\begin{proof}
$\|m_q\|_{\HS}^2 = \sum_{a,b,c} q^{a+b+c}
\|m_{a,b}^c\|_{\HS}^2 = \Sigma_{\HS}(q) < \infty$.
A genus-$g$ surface with $n$ punctures decomposes into
$2g - 2 + n$ pairs of pants; each sewing circle contributes
one composition of sector maps.  Composition of two
Hilbert--Schmidt operators is trace class:
$\|AB\|_1 \leq \|A\|_{\HS} \|B\|_{\HS}$.  Inductively,
sewing $k$ circles produces an operator of trace norm at most
$\Sigma_{\HS}(q)^{k/2}$.  For a closed surface ($n = 0$),
$k = 3g - 3$, and the amplitude is
$\Tr(\text{trace-class}) < \infty$.
\end{proof}

\begin{theorem}[General HS-sewing criterion;
\ClaimStatusProvedElsewhere]%
\label{thm:g1-general-hs}%
\index{HS-sewing!general criterion}%
Let\/ $\cA$ be a positive-energy chiral algebra with
\begin{enumerate}[label=\textup{(\roman*)}]
\item subexponential sector growth:
  $\log\dim H_n = o(n)$;
\item polynomial OPE growth:
  $|C^{c,k}_{a,i;\,b,j}| \leq K(a+b+c+1)^N$ for constants
  $K, N$.
\end{enumerate}
Then $\cA$ satisfies HS-sewing for every $0 < q < 1$.
\end{theorem}

\begin{proof}
Theorem~\ref{thm:general-hs-sewing}.  The bound is:
\begin{align}\label{eq:g1-hs-bound}
\|m_{a,b}^c\|_{\HS}^2
&\;\leq\;
\dim H_a \cdot \dim H_b \cdot \dim H_c
\cdot K^2(a{+}b{+}c{+}1)^{2N}, \\
\Sigma_{\HS}(q)
&\;\leq\;
K^2 \Bigl(\sum_n q^n \dim H_n (n{+}1)^{2N}\Bigr)
\Bigl(\sum_n q^n \dim H_n\Bigr)^2 \;<\; \infty.
\end{align}
Hypothesis (i) gives $\dim H_n \leq Ce^{\alpha\sqrt{n}}$;
$q^n$ decays exponentially; each factor converges for all
$q < 1$.
\end{proof}

\begin{definition}[Sewing operator at genus $g$]%
\label{def:g1-sewing-operator}%
\index{sewing operator|textbf}%
Fix a decomposition of a closed genus-$g$ surface $\Sigma_g$ into
$2g - 2$ pairs of pants $P_1, \ldots, P_{2g-2}$ with sewing
circles $\gamma_1, \ldots, \gamma_{3g-3}$.  For each
circle~$\gamma_j$ with collar parameter $q_j$, define the
\emph{sewing operator}
\begin{equation}\label{eq:g1-sewing-op}
T_j \;:=\; \sum_{n \geq 0} q_j^n\,\pi_{n},
\qquad
\|T_j\|_1 = \sum_{n \geq 0} q_j^n \dim H_n,
\end{equation}
where $\pi_n \colon H \to H_n$ is the projection onto weight
$n$.  Under HS-sewing, $T_j$ is trace class.
\end{definition}

\begin{definition}[Gravitational amplitude at order $(r,g)$]%
\label{def:g1-grav-amplitude}%
\index{gravitational amplitude|textbf}%
\index{graph amplitude!gravitational}%
For a Koszul chiral algebra $\cA$ with universal MC element
$\Theta_\cA \in \MC(\gAmod)$
(Theorem~\ref{thm:mc2-bar-intrinsic}), the
\emph{gravitational amplitude at order $(r,g)$} is
\begin{equation}\label{eq:g1-grav-amplitude}
\mathcal{Z}_{r,g}(\cA)
\;:=\;
\sum_{\substack{\Gamma \in \mathcal{G}_{g,n}^{\mathrm{st}}
\\ w(\Gamma) = r}}
\frac{\hbar^{g(\Gamma)}}{|\Aut(\Gamma)|}
\;\mathcal{A}_\Gamma
\;\in\;
\cA^{\mathrm{sh}}_{r,g},
\end{equation}
where $\mathcal{A}_\Gamma$ is the graph amplitude of~$\Gamma$:
the integral over the FM compactification
$\overline{C}_{|V(\Gamma)|}(X)$ of the product of propagators
$\prod_{e \in E(\Gamma)} \eta_e$ contracted with vertex labels
$\prod_{v \in V(\Gamma)} m_{\operatorname{val}(v)}^{(g(v))}$
from the transferred cyclic $\Ainf$ operations
(Definition~\ref{def:modular-bar-hamiltonian}).  The weight
$w(\Gamma) = |V(\Gamma)| + 2b_1(\Gamma)$ controls the filtration.
\end{equation}
\end{definition}

\begin{proposition}[Logarithmic integrability of graph
amplitudes; \ClaimStatusProvedHere]%
\label{prop:g1-log-integrability}%
\index{graph amplitude!logarithmic integrability|textbf}%
\index{UV finiteness!logarithmic integrability}%
For any stable graph
$\Gamma \in \mathcal{G}_{g,n}^{\mathrm{st}}$ and any smooth
projective curve $X$, the graph amplitude $\mathcal{A}_\Gamma$
converges absolutely.
\end{proposition}

\begin{proof}
The integrand is a section of
$\Omega^{|E(\Gamma)|}(\log D) \otimes \mathcal{V}$ on
$\overline{C}_{|V|}(X)$, where $D$ is the normal-crossing
boundary (Theorem~\ref{thm:FM}(3)) and $\mathcal{V}$ is the
algebraic coefficient sheaf from the vertex labels.

\emph{Local analysis near a collision stratum.}
Let $D_S$ be the boundary divisor for a cluster $S \subset V$
with $|S| \geq 2$.  In the local coordinates of
Theorem~\ref{thm:FM}(5)---collar parameter $\epsilon$,
angular coordinates $\theta$ on
$\overline{C}_{|S|}(\mathbb{A}^1)$---the propagator
$\eta_e$ for an edge $e$ within the cluster has the form
\begin{equation}\label{eq:g1-local-eta}
\eta_e
\;=\;
\frac{d\epsilon}{\epsilon}
\;+\;
\omega_{\mathrm{ang}}
\;+\;
O(\epsilon),
\end{equation}
where $\omega_{\mathrm{ang}}$ is smooth on
$\overline{C}_{|S|}(\mathbb{A}^1)$.  The leading pole is simple
($d\epsilon/\epsilon$), and $D$ has normal crossings, so each
$\eta_e$ has at most a logarithmic singularity.

\emph{Arnold cancellation.}
The Arnold relation $\eta_{ij} \wedge \eta_{jk} +
\eta_{jk} \wedge \eta_{ki} + \eta_{ki} \wedge \eta_{ij} = 0$
(Theorem~\ref{thm:arnold-relations}, at genus~$0$; corrected by
$\kappa \cdot \omega_g$ at genus~$g \geq 1$) forces cancellations
among the poles of the product $\bigwedge_e \eta_e$.  At a
codimension-$k$ stratum, the total pole order of the integrand
is at most $k$ (matching the codimension), so the
integral has at most a logarithmic divergence---which is
integrable on a compact manifold with normal-crossing boundary
(Deligne's $L^2$ integrability of logarithmic forms).

\emph{Global finiteness.}
$\overline{C}_{|V|}(X)$ is smooth and proper
(Theorem~\ref{thm:FM}(2)).  A section of $\Omega^p(\log D)$ on
a smooth proper variety with normal-crossing boundary~$D$ is
$L^1$-integrable.  The vertex labels are bounded algebraic data.
The $L^1$-norm is finite; the integral converges absolutely.
\end{proof}

\begin{proposition}[Arity cutoff; \ClaimStatusProvedElsewhere]%
\label{prop:g1-arity-cutoff}%
\index{arity cutoff!perturbative finiteness}%
For a strong completion tower $\cA = \varprojlim_N \cA_{\leq N}$
\textup{(}Definition~\textup{\ref{def:strong-completion-tower})},
the MC equation on $\cA_{\leq N}$ involves only arities
$r \leq N$
\textup{(}Lemma~\textup{\ref{lem:arity-cutoff})}:
the graph sum~\eqref{eq:g1-grav-amplitude} restricted to
$\cA_{\leq N}$ has $w(\Gamma) \leq N$.
At each finite stage, the perturbative amplitude is a finite
sum of convergent integrals.
\end{proposition}

\begin{proof}
The strong filtration axiom
$\mu_r(F^{i_1}, \ldots, F^{i_r}) \subset
F^{i_1 + \cdots + i_r}$
ensures that modulo $F^{N+1}$, all terms with
$\sum_v \operatorname{val}(v) > N$ vanish.  The number of
graphs at weight $\leq N$ is finite (bounded by $(2N)^{2N}$).
\end{proof}

% ======================================================================
%
%  STAGE 2: SHADOW FREE ENERGIES AND THE BERNOULLI
%           GENERATING FUNCTION
%
% ======================================================================

\subsection{Shadow free energies and the Bernoulli generating
function}
\label{subsec:thqg-I-shadow-free-energies}
\index{shadow free energy|textbf}
\index{Bernoulli generating function|textbf}

The scalar shadow of the MC element $\Theta_\cA$ at genus $g$
is the modular characteristic $\kappa(\cA) \cdot \lambda_g$
(Theorem~\ref{thm:genus-universality}).  The genus generating
function of these scalar shadows is controlled by the
$\hat{A}$-genus.

\begin{definition}[Shadow free energy]%
\label{def:g1-shadow-free-energy}%
\index{shadow free energy!definition}%
The \emph{shadow free energy} of a Koszul chiral algebra
$\cA$ is the formal power series
\begin{equation}\label{eq:g1-shadow-free-energy}
\mathcal{F}^{\mathrm{sh}}(\cA;\hbar)
\;:=\;
\sum_{g \geq 1} \hbar^g \cdot
\Tr\bigl(\Theta_\cA^{(g)}\bigr)
\;=\;
\sum_{g \geq 1} \hbar^g \cdot \kappa(\cA) \cdot
\int_{\overline{\mathcal{M}}_g} \lambda_g
\;\in\;
\hbar\,\mathbb{Q}[\![\hbar]\!].
\end{equation}
The trace extracts the scalar component from the
$\gAmod$-valued MC element.
\end{definition}

\begin{theorem}[$\hat{A}$-genus generating function;
\ClaimStatusProvedElsewhere]%
\label{thm:g1-ahat-generating}%
\index{$\hat{A}$-genus!generating function}%
\index{Hodge integral!$\lambda_g$ class}%
The generating function of Hodge integrals
$\int_{\overline{\mathcal{M}}_g} \lambda_g$ is
\begin{equation}\label{eq:g1-ahat-gf}
\sum_{g \geq 1} \hbar^g
\int_{\overline{\mathcal{M}}_g} \lambda_g
\;=\;
\log\!\Bigl(
\frac{\hbar/2}{\sinh(\hbar/2)}
\Bigr)
\;=\;
\sum_{g \geq 1}
\frac{(-1)^{g+1} B_{2g}}{2g(2g)!}\,\hbar^{2g},
\end{equation}
where $B_{2g}$ are the Bernoulli numbers:
$B_2 = \tfrac{1}{6}$, $B_4 = -\tfrac{1}{30}$,
$B_6 = \tfrac{1}{42}$, $B_8 = -\tfrac{1}{30}$.
\end{theorem}

\begin{proof}
Theorem~\ref{thm:family-index}:
$\sum_{g \geq 0} \kappa(\cA)\lambda_g =
\kappa(\cA)(\hat{A}(i\hbar) - 1)$,
where $\hat{A}(x) = (x/2)/\sinh(x/2)$. Taking logarithm and
evaluating $\int_{\overline{\mathcal{M}}_g}$ gives the
Bernoulli series.  The Hodge integral
$\int_{\overline{\mathcal{M}}_g} \lambda_g$ is
Faber--Pandharipande's evaluation~\cite{FP00}:
$\int_{\overline{\mathcal{M}}_g} \lambda_g =
(-1)^{g+1} B_{2g}/(2g(2g)!)$ for $g \geq 2$, and
$\int_{\overline{\mathcal{M}}_1} \lambda_1 = 1/24$.
\end{proof}

\begin{corollary}[Explicit shadow free energy;
\ClaimStatusProvedHere]%
\label{cor:g1-explicit-free-energy}%
\index{shadow free energy!explicit values}%
\begin{equation}\label{eq:g1-explicit-free-energy}
\mathcal{F}^{\mathrm{sh}}(\cA;\hbar)
\;=\;
\kappa(\cA) \cdot
\log\!\Bigl(
\frac{\hbar/2}{\sinh(\hbar/2)}
\Bigr).
\end{equation}
The first four terms:
\begin{equation}\label{eq:g1-free-energy-terms}
\mathcal{F}^{\mathrm{sh}}
\;=\;
\kappa\!\left[
\frac{\hbar^2}{24}
- \frac{\hbar^4}{2880}
+ \frac{\hbar^6}{181440}
- \frac{\hbar^8}{9676800}
+ \cdots
\right].
\end{equation}
\end{corollary}

\begin{proof}
Substitute $\int_{\overline{\mathcal{M}}_g}\lambda_g$
from Theorem~\ref{thm:g1-ahat-generating} into
Definition~\ref{def:g1-shadow-free-energy}.
Genus~$1$: $\kappa \cdot 1/24 = \kappa \hbar^2/24$.
Genus~$2$: $\kappa \cdot (-B_4)/(4 \cdot 24) =
\kappa \cdot 1/(4 \cdot 24 \cdot 30) = \kappa/(2880)$
with sign $(-1)^3 = -1$, giving $-\kappa\hbar^4/2880$.
\end{proof}

\begin{definition}[Shadow partition function]%
\label{def:g1-shadow-partition}%
\index{shadow partition function|textbf}%
The \emph{shadow partition function} is the
exponentiated free energy:
\begin{equation}\label{eq:g1-shadow-partition}
Z^{\mathrm{sh}}(\cA;\hbar)
\;:=\;
\exp\!\bigl(\mathcal{F}^{\mathrm{sh}}(\cA;\hbar)\bigr)
\;=\;
\Bigl(\frac{\hbar/2}{\sinh(\hbar/2)}\Bigr)^{\kappa(\cA)}.
\end{equation}
\end{definition}

\begin{remark}[The shadow partition function as
$\hat{A}$-power]%
\label{rem:g1-ahat-power}%
\index{$\hat{A}$-genus!partition function}%
$Z^{\mathrm{sh}}(\cA;\hbar) = \hat{A}(i\hbar)^{\kappa(\cA)}$.
The modular characteristic $\kappa(\cA)$ is the exponent of the
universal multiplicative genus.  This is the content of
Theorem~D: the modular characteristic is universal, additive,
and determines the $\hat{A}$-class generating function.
\end{remark}

\begin{proposition}[Genus-$g$ scalar shadow: explicit values;
\ClaimStatusProvedHere]%
\label{prop:g1-scalar-shadow-values}%
\index{modular characteristic!genus-$g$ values}%
\begin{center}
\renewcommand{\arraystretch}{1.4}
\begin{tabular}{c|c|c}
$g$ & $\int_{\overline{\mathcal{M}}_g} \lambda_g$ &
  $\Tr(\Theta_\cA^{(g)})$
\\[2pt]\hline\hline
$1$ & $1/24$ & $\kappa/24$ \\
$2$ & $-1/2880$ & $-\kappa/2880$ \\
$3$ & $1/181440$ & $\kappa/181440$ \\
$4$ & $-1/9676800$ & $-\kappa/9676800$ \\
$5$ & $1/479001600$ & $\kappa/479001600$
\end{tabular}
\end{center}
These values satisfy $(-1)^{g+1} > 0$: the
sign alternates, with the genus-$1$ contribution positive.
\end{proposition}

\begin{proof}
$\int_{\overline{\mathcal{M}}_g} \lambda_g =
(-1)^{g+1} B_{2g}/(2g(2g)!)$.  The Bernoulli numbers are
$B_2 = 1/6$, $B_4 = -1/30$, $B_6 = 1/42$,
$B_8 = -1/30$, $B_{10} = 5/66$.  Substituting:
$g = 1$: $B_2/(2 \cdot 2!) = (1/6)/4 = 1/24$.
$g = 2$: $-B_4/(4 \cdot 4!) = (1/30)/96 = 1/2880$,
with sign $(-1)^3 = -1$.  And so on.
\end{proof}

\begin{proposition}[Nonlinear corrections to the shadow free
energy; \ClaimStatusProvedHere]%
\label{prop:g1-nonlinear-corrections}%
\index{shadow free energy!nonlinear corrections}%
\index{shadow Postnikov tower!free energy corrections}%
Beyond the scalar shadow, the shadow Postnikov tower contributes
nonlinear corrections.  The full free energy is
\begin{equation}\label{eq:g1-full-free-energy}
\mathcal{F}(\cA;\hbar)
\;=\;
\mathcal{F}^{\mathrm{sh}}(\cA;\hbar)
\;+\;
\sum_{r \geq 3}
\mathcal{F}^{(r)}(\cA;\hbar),
\end{equation}
where $\mathcal{F}^{(r)}$ collects contributions from graphs of
weight exactly~$r$.  The classification:
\begin{center}
\renewcommand{\arraystretch}{1.3}
\begin{tabular}{c|c|c}
Shadow depth class & $r_{\max}$ & Nonlinear corrections
\\[2pt]\hline\hline
G \textup{(}Gaussian\textup{)} & $2$ &
  $\mathcal{F}^{(r)} = 0$ for $r \geq 3$ \\
L \textup{(}Lie/tree\textup{)} & $3$ &
  $\mathcal{F}^{(3)} \neq 0$,
  $\mathcal{F}^{(r)} = 0$ for $r \geq 4$ \\
C \textup{(}contact\textup{)} & $4$ &
  $\mathcal{F}^{(3)}, \mathcal{F}^{(4)} \neq 0$,
  $\mathcal{F}^{(r)} = 0$ for $r \geq 5$ \\
M \textup{(}mixed\textup{)} & $\infty$ &
  $\mathcal{F}^{(r)} \neq 0$ for all $r$
\end{tabular}
\end{center}
\end{proposition}

\begin{proof}
The shadow tower
$\Theta_\cA = \Theta_\cA^{\leq 2} + \Theta_\cA^{(3)} +
\Theta_\cA^{(4)} + \cdots$ terminates at $r_{\max}$
(Definition~\ref{def:shadow-postnikov-tower}).  The
free-energy correction at arity~$r$ is
$\mathcal{F}^{(r)} = \sum_g \hbar^g \Tr(\Theta_\cA^{(r),(g)})$,
which vanishes identically when the shadow tower terminates
below arity~$r$.  The classification follows from the shadow
depth values: $r_{\max} = 2$ for Heisenberg/lattice (G class),
$3$ for affine Kac--Moody (L class), $4$ for $\beta\gamma$
(C class), $\infty$ for Virasoro and $\mathcal{W}_N$ (M class).
\end{proof}

% ======================================================================
%
%  STAGE 3: CONVERGENCE OF THE GRAVITATIONAL PARTITION
%           FUNCTION
%
% ======================================================================

\subsection{Convergence of the gravitational partition function}
\label{subsec:thqg-I-convergence-grav}
\index{partition function!convergence|textbf}
\index{perturbative finiteness!convergence}

\begin{theorem}[Perturbative finiteness; \ClaimStatusProvedHere]%
\label{thm:g1-perturbative-finiteness}%
\index{perturbative finiteness!main theorem|textbf}%
Let\/ $\cA$ be a chiral algebra on a smooth projective curve $X$
satisfying:
\begin{enumerate}[label=\textup{(\roman*)}]
\item \emph{Koszulness:} $(\cA, \cA^!)$ is a chiral Koszul pair
  \textup{(}Definition~\textup{\ref{def:koszul-chiral-pair})};
\item \emph{HS-sewing:}
  $\Sigma_{\HS}(q) < \infty$ for some $0 < q < 1$
  \textup{(}Definition~\textup{\ref{def:g1-hs-sewing})}.
\end{enumerate}
Then:
\begin{enumerate}[label=\textup{(\alph*)}]
\item Each graph amplitude $\mathcal{A}_\Gamma$ converges
  absolutely
  \textup{(}Proposition~\textup{\ref{prop:g1-log-integrability})}.
\item At fixed arity $r$ and genus $g$, the amplitude
  $\mathcal{Z}_{r,g}(\cA)$ is a finite sum of convergent
  integrals.
\item The arity series
  $\mathcal{Z}_g(\cA) := \sum_{r \geq 2}
  \mathcal{Z}_{r,g}(\cA)$ converges.
\item The genus series
  $Z(\cA;\hbar) := \sum_{g \geq 0} \hbar^g Z_g(\cA)$
  converges for $|\hbar|$ sufficiently small.
\item No UV renormalization is required.
\end{enumerate}
\end{theorem}

\begin{proof}
\emph{(a)} Proposition~\ref{prop:g1-log-integrability}.

\emph{(b)} At weight $r$, the number of contributing stable
graphs satisfies $|\{\Gamma : w(\Gamma) = r\}| \leq (2r)^{2r}$
(the graph count: a connected graph on $v \leq r$ labeled
vertices with $e \leq r$ edges has at most $r^{2r}$ choices;
genus labels and half-edges contribute at most a further
$(2r)^r$ factor).  A finite sum of convergent integrals converges.

\emph{(c)} The individual amplitude bound
(Proposition~\ref{prop:g1-individual-bound} below) gives
$|\mathcal{A}_\Gamma| \leq C^{|V|} q^{|E|/2}$ with the HS-sewing
constant $C = \Sigma_{\HS}(q)^{1/2}$ and collar parameter $q$.
The arity-$r$ bound:
\begin{equation}\label{eq:g1-arity-bound}
|\mathcal{Z}_{r,g}|
\;\leq\;
(2r)^{2r} C^r q^{r/4},
\end{equation}
since $|E| \geq |V| - 1 \geq r/2 - 1$.  For fixed $q < 1$,
the super-polynomial $(2r)^{2r}$ is dominated by $q^{r/4}$
(exponential decay): $(2r)^{2r} q^{r/4} = \exp(2r\log(2r) -
(r/4)\log(1/q)) \to 0$ since $\log(2r)/r \to 0$.

\emph{(d)} The genus-$g$ amplitude has an additional factor from
the propagator's genus dependence.  On $\Sigma_g$, the
propagator satisfies $|\eta^{(g)}(z,w)| \leq C_g/|z-w| + C_g'$
with $C_g = O(g)$ (Proposition~\ref{prop:g1-propagator-bound}).
The graph bound at genus $g$ becomes
$|\mathcal{Z}_{r,g}| \leq (2r)^{2r} C^r g^{O(r)} q^{r/4}$.
Summing over $r$ and $g$:
$\sum_{g,r} |\hbar|^g |\mathcal{Z}_{r,g}| \leq
\sum_r (2r)^{2r} C^r q^{r/4}
\sum_g |\hbar|^g g^{O(r)} < \infty$
provided $|\hbar| < 1$ (the genus sum is polynomial in $g$ at
each $r$, so converges for $|\hbar|$ sufficiently small).

\emph{(e)} The FM compactification provides the regulator
(Theorem~\ref{thm:FM}).  No counterterms, subtractions, or
renormalization group actions appear in the proof.
\end{proof}

\begin{proposition}[Individual graph amplitude bound;
\ClaimStatusProvedHere]%
\label{prop:g1-individual-bound}%
\index{graph amplitude!individual bound}%
Let $\Gamma$ be a stable graph with $v = |V|$ vertices,
$e = |E|$ edges.  Under HS-sewing at collar parameter $q$,
\begin{equation}\label{eq:g1-individual-bound}
|\mathcal{A}_\Gamma|
\;\leq\;
\frac{v! \cdot (4\pi)^v}{|\Aut(\Gamma)|}
\;\cdot\;
\|\eta\|_q^e
\;\cdot\;
\prod_{w \in V} \|m_{\operatorname{val}(w)}^{(g(w))}\|_q,
\end{equation}
where $\|\eta\|_q$ is the propagator norm and
$\|m_k^{(h)}\|_q$ is the vertex norm in the $q$-weighted
topology.
\end{proposition}

\begin{proof}
Factor the graph amplitude into propagator and vertex
contributions.  The integral over $\overline{C}_v(X)$ of
the angular part is bounded by $\Vol(\overline{C}_v(X))
\leq v! (4\pi)^v$ (the volume of the FM compactification
on a curve of area~$1$; each blow-up increases volume by at
most $4\pi$).  The propagator contributes $\|\eta\|_q$ per
edge; the vertices contribute $\|m_k^{(h)}\|_q$ per vertex.
\end{proof}

\begin{proposition}[Propagator bound at genus $g$;
\ClaimStatusProvedHere]%
\label{prop:g1-propagator-bound}%
\index{propagator!genus-$g$ bound}%
On $\Sigma_g$ with period matrix $\Omega$ and Arakelov metric,
the single-valued propagator satisfies
\begin{equation}\label{eq:g1-propagator-bound}
|\eta^{(g)}(z,w)|
\;\leq\;
\frac{1 + \pi g\,\|\operatorname{Im}\Omega^{-1}\|}
{|z - w|}
\;+\;
\pi g\,\|\operatorname{Im}\Omega^{-1}\|.
\end{equation}
\end{proposition}

\begin{proof}
The propagator is $\eta^{(g)}_{ij} =
[\partial_{z_i}\log E(z_i,z_j) +
\pi\sum_{\alpha,\beta}
\omega_\alpha(z_i)(\operatorname{Im}\Omega)^{-1}_{\alpha\beta}
\operatorname{Im}(\int_{z_0}^{z_j}\omega_\beta)]
(dz_i - dz_j)$.
The prime form satisfies $|E(z,w)| \sim |z - w|$ near the
diagonal, giving $|\partial_z \log E| \leq 1/|z-w| + O(1)$.
The correction term is bounded by
$\pi \sum_{\alpha,\beta} |\omega_\alpha|
|(\operatorname{Im}\Omega)^{-1}_{\alpha\beta}|
|\int \omega_\beta| \leq
\pi g \|\operatorname{Im}\Omega^{-1}\|$ (using
$|\omega_\alpha| \leq 1$ in the Arakelov metric and
$|\int_{z_0}^{z_j} \omega_\beta| \leq 1$ on the
fundamental domain).
\end{proof}

\begin{corollary}[Genus-$g$ total amplitude bound;
\ClaimStatusProvedHere]%
\label{cor:g1-genus-bound}%
\index{genus-$g$!amplitude bound}%
\begin{equation}\label{eq:g1-genus-total-bound}
\sum_{r \geq 2}
|\mathcal{Z}_{r,g}(\cA)|
\;\leq\;
\sum_{r \geq 2}
(2r)^{2r}\,C^r\,g^{O(r)}\,q^{r/4}
\;<\;\infty
\end{equation}
for each $g$ and sufficiently small $q$.
The genus series $Z(\cA;\hbar)$ converges for
$|\hbar| < \hbar_0(q)$.
\end{corollary}

\begin{proof}
At genus $g$, the propagator norm is $O(g)$
(Proposition~\ref{prop:g1-propagator-bound}).  In the graph
bound, each edge contributes $O(g)$ instead of $O(1)$, giving
an extra factor $g^{|E|} \leq g^{O(r)}$ at weight $r$.
The exponential decay $q^{r/4}$ still dominates for fixed $g$.
For the genus sum: $\sum_g |\hbar|^g g^{O(r)}$ converges for
$|\hbar| < 1$ (polynomial growth in $g$).
\end{proof}

\begin{theorem}[Completed perturbative finiteness;
\ClaimStatusProvedHere]%
\label{thm:g1-completed-finiteness}%
\index{perturbative finiteness!completed tower|textbf}%
For a strong completion tower
$\cA = \varprojlim_N \cA_{\leq N}$
\textup{(}Definition~\textup{\ref{def:strong-completion-tower})}
with each $\cA_{\leq N}$ Koszul and HS-sewing:
\begin{enumerate}[label=\textup{(\alph*)}]
\item The completed MC element
  $\widehat{\Theta}_\cA = \varprojlim_N
  \Theta_{\cA_{\leq N}}$ exists
  \textup{(}Theorem~\textup{\ref{thm:completed-bar-cobar-strong})}.
\item The completed amplitudes
  $\widehat{\mathcal{Z}}_{r,g}(\cA) =
  \mathcal{Z}_{r,g}(\cA_{\leq r})$
  are finite for all $(r,g)$, by coefficient stabilization:
  $\mathcal{Z}_{r,g}(\cA_{\leq N})$ is independent of $N$
  for $N \geq r$.
\item The MC equation holds on the completed amplitudes.
\end{enumerate}
\end{theorem}

\begin{proof}
Part (a): Theorem~\ref{thm:completed-bar-cobar-strong}.
Part (b): by the arity cutoff
(Proposition~\ref{prop:g1-arity-cutoff}), the weight-$r$
amplitude depends only on vertex labels of valence $\leq r$,
which are determined by $\cA_{\leq r}$.  For $N \geq r$, the
projection $\cA_{\leq N} \to \cA_{\leq r}$ does not affect
arities $\leq r$.  Part (c): the MC equation holds at each
stage and passes to the inverse limit (completeness of the
convolution algebra).
\end{proof}

% ======================================================================
%
%  STAGE 4: THE HEISENBERG PARTITION FUNCTION AS PROTOTYPE
%
% ======================================================================

\subsection{The Heisenberg partition function as prototype}
\label{subsec:thqg-I-heisenberg-prototype}
\index{Heisenberg!partition function|textbf}
\index{Gaussian shadow!partition function}
\index{Fredholm determinant!Heisenberg}

The Heisenberg algebra $\mathcal{H}_k$ is the Gaussian
prototype: shadow depth $r_{\max} = 2$, all higher operations
vanish, and the full computation reduces to second quantization.

\begin{proposition}[Heisenberg HS-sewing;
\ClaimStatusProvedElsewhere]%
\label{prop:g1-heisenberg-hs}%
\index{Heisenberg!HS-sewing}%
$\mathcal{H}_k$ satisfies HS-sewing at all $q$ with
$\Sigma_{\HS}(q) = k^2/(1-q)^3$.  The sector maps are
$m_{a,b}^c = k\,\delta_{a+b,c}$ (free-field Wick contraction),
so $\|m_{a,b}^c\|_{\HS}^2 = k^2\,p(a)\,p(b)\,\delta_{a+b,c}$
and
\begin{equation}\label{eq:g1-heisenberg-hs}
\Sigma_{\HS}(q)
\;=\;
k^2 \sum_{a,b \geq 0} q^{2a+2b}\,p(a)\,p(b)
\;=\;
k^2 \Bigl(\sum_n q^{2n} p(n)\Bigr)^2
\;=\;
\frac{k^2}{\prod_{n \geq 1}(1-q^{2n})^2}.
\end{equation}
\end{proposition}

\begin{proposition}[Heisenberg shadow tower terminates;
\ClaimStatusProvedElsewhere]%
\label{prop:g1-heisenberg-terminates}%
\index{Heisenberg!shadow termination}%
The transferred cyclic $\Ainf$ operations of $\mathcal{H}_k$
satisfy $m_k^{(h)} = 0$ for $k \geq 3$.
$\Theta_{\mathcal{H}_k} = \Theta^{\leq 2} =
k \cdot \eta \otimes \Lambda$
with $\kappa(\mathcal{H}_k) = k$.  The graph sum reduces to
the single self-loop graph $\Gamma_{\mathrm{irr}}$ and its
higher-genus iterates:
\begin{equation}\label{eq:g1-heisenberg-theta}
\Theta_{\mathcal{H}_k}^{(g)}
\;=\;
k \cdot \eta \otimes \lambda_g,
\qquad
\Tr(\Theta^{(g)}) = k/24 \;\text{at}\; g = 1.
\end{equation}
\end{proposition}

\begin{theorem}[Heisenberg one-particle sewing;
\ClaimStatusProvedElsewhere]%
\label{thm:g1-heisenberg-sewing}%
\index{Heisenberg!one-particle sewing|textbf}%
\index{Bergman space!one-particle}%
The pair-of-pants amplitude of $\mathcal{H}_k$ is the
second quantization $\Gamma(R)$ of the one-particle Bergman
restriction $R \colon A^2(D_{\mathrm{out}}) \to
A^2(D_{\mathrm{in},1}) \otimes A^2(D_{\mathrm{in},2})$,
where $A^2(D) = \{f \text{ holomorphic on } D :
\int_D |f|^2 < \infty\}$ is the Bergman space of a disk~$D$.
The one-particle amplitudes satisfy
$|R_{n,m}| \leq Cq^{(n+m)/2}$.  The sewing operator $T_q$ on
$A^2(\partial D)$ has eigenvalues $q^n$, $n \geq 1$, and is
trace class: $\|T_q\|_1 = \sum_n q^n = q/(1-q)$.
\end{theorem}

\begin{proof}
Theorem~\ref{thm:heisenberg-one-particle-sewing}.
Wick's theorem: all amplitudes factor into pairings of the free
field $J(z)$.  The one-particle restriction is computed by the
Bergman kernel:
$R_{n,m} = \frac{k}{2\pi i}\oint \oint z^{-n-1}w^{-m-1}
G_{\mathrm{pants}}(z,w)\,dz\,dw$.
The Green function on the pair of pants satisfies
$|G_{\mathrm{pants}}(z,w)| \leq C|z|^{-1/2}|w|^{-1/2}$ in the
collar region, giving $|R_{n,m}| \leq Cq^{(n+m)/2}$.
\end{proof}

\begin{theorem}[Heisenberg genus-$g$ partition function;
\ClaimStatusProvedElsewhere]%
\label{thm:g1-heisenberg-zg}%
\index{Heisenberg!genus-$g$ partition function|textbf}%
\index{Fredholm determinant!genus-$g$}%
\begin{equation}\label{eq:g1-heisenberg-zg}
Z_g(\mathcal{H}_k;\,\Omega)
\;=\;
(\det\operatorname{Im}\Omega)^{-k/2}
\;\cdot\;
(\det{}'_\zeta \Delta_{\Sigma_g})^{-k/2},
\end{equation}
where $\Delta_{\Sigma_g}$ is the scalar Laplacian on $\Sigma_g$
and $\det{}'_\zeta$ the zeta-regularized determinant with zero
mode removed.
\end{theorem}

\begin{proof}
The genus-$g$ amplitude sews $2g - 2$ pair-of-pants operators
along $3g - 3$ circles.  Each sewing is a trace over the Fock
space: $\Tr_{\mathcal{F}}(\Gamma(T_q)) =
\Det(1 - T_q)^{-k}$.  For $g$ sewings with sewing operators
$T_1, \ldots, T_{3g-3}$, the partition function is the
Fredholm determinant
$Z_g = \prod_j \Det(1 - T_j)^{-k/2}$.  By the Arakelov
factorization (the period matrix determines all sewing
parameters), this equals
$(\det\operatorname{Im}\Omega)^{-k/2}
(\det{}'_\zeta\Delta)^{-k/2}$.
\end{proof}

\begin{corollary}[Heisenberg at genus $1$;
\ClaimStatusProvedElsewhere]%
\label{cor:g1-heisenberg-genus1}%
\index{Heisenberg!genus-$1$ partition function}%
\begin{equation}\label{eq:g1-heisenberg-z1}
Z_1(\mathcal{H}_k;\tau)
\;=\;
(\operatorname{Im}\tau)^{-k/2}
\,|\eta(\tau)|^{-2k},
\qquad
\eta(\tau) = q^{1/24}\prod_{n=1}^\infty(1-q^n).
\end{equation}
\end{corollary}

\begin{proof}
At genus~$1$, $\Omega = \tau \in \mathfrak{H}_1$,
$\det\operatorname{Im}\Omega = \operatorname{Im}\tau$, and
$\det'_\zeta\Delta_{E_\tau} = |\eta(\tau)|^4$ (by the
Kronecker limit formula).
\end{proof}

\begin{remark}[Comparison with the shadow free energy]%
\label{rem:g1-heisenberg-comparison}%
\index{Heisenberg!shadow comparison}%
The shadow partition function
$Z^{\mathrm{sh}}(\mathcal{H}_k;\hbar) =
(\hbar/2/\sinh(\hbar/2))^k$ captures the scalar part.
The full partition function $Z_g(\mathcal{H}_k;\Omega)$
depends on the period matrix $\Omega$; the shadow partition
function depends only on $\hbar$ and $\kappa = k$.
The relationship: integrating
$Z_g(\mathcal{H}_k;\Omega)$ over
$\overline{\mathcal{M}}_g$ with the Weil--Petersson measure
recovers $Z^{\mathrm{sh}}_g(\mathcal{H}_k)$ at the level
of tautological classes.
\end{remark}

% ======================================================================
%
%  STAGE 5: FINITENESS FOR THE STANDARD LANDSCAPE
%
% ======================================================================

\subsection{Finiteness for the standard landscape}
\label{subsec:thqg-I-standard-landscape}
\index{perturbative finiteness!standard landscape|textbf}
\index{standard landscape!finiteness}

\begin{theorem}[Finiteness for the standard landscape;
\ClaimStatusProvedHere]%
\label{thm:g1-finiteness-landscape}%
\index{perturbative finiteness!five families|textbf}%
Every algebra in the standard landscape satisfies the hypotheses
of Theorem~\textup{\ref{thm:g1-perturbative-finiteness}}:
\begin{center}
\renewcommand{\arraystretch}{1.4}
\begin{tabular}{c|c|c|c|c|c|c}
$\cA$ & $\kappa(\cA)$ & $r_{\max}$ & $\dim H_n$
  & $|C_{a,b}^c|$ & HS & Koszul
\\[2pt]\hline\hline
$\mathcal{H}_k$ & $k$ & $2$
  & $p(n)$ & $\leq k$ & $\checkmark$ & $\checkmark$
\\
$V_k(\fg)$ &
  $\frac{(k{+}h^\vee)\dim\fg}{2h^\vee}$
  & $3$ & $\leq Ce^{\alpha\sqrt{n}}$
  & $\leq K(n{+}1)^{h^\vee}$ & $\checkmark$ & $\checkmark$
\\
$\mathrm{Vir}_c$ & $\frac{5c}{6}$
  & $\infty$ & $p(n)$
  & $\leq K(n{+}1)^4$ & $\checkmark$ & $\checkmark$
\\
$\mathcal{W}_N$ &
  $\frac{(N{-}1)(N^2{+}N{+}k(N^2{-}1))}{N+k}$
  & $\infty$ & $\leq Ce^{\alpha\sqrt{n}}$
  & $\leq K(n{+}1)^{2N}$ & $\checkmark$ & $\checkmark$
\\
$V_\Lambda$ & $\rank(\Lambda)$ & $2$
  & $|\Lambda^*/\Lambda|\,p(n)^r$
  & $\leq kr$ & $\checkmark$ & $\checkmark$
\end{tabular}
\end{center}
Consequently, all genus-$g$ amplitudes in every family are finite
at every order in the shadow tower.
\end{theorem}

\begin{proof}
\emph{HS-sewing}: Theorem~\ref{thm:g1-general-hs} applies to
each family.  All five have subexponential sector growth
($p(n) \leq Ce^{\pi\sqrt{2n/3}}$ for the partition function;
$\dim H_n \leq Ce^{\alpha\sqrt{n}}$ for Kac--Moody and
$\mathcal{W}$-algebras by the character formulas) and polynomial
OPE growth (the OPE of generators has finitely many singular
terms; the growth exponents are recorded in the table).

\emph{Koszulness}: $\mathcal{H}_k$ by
Theorem~\ref{thm:km-chiral-koszul} (free algebra, diagonal Ext
vanishing is immediate).
$V_k(\fg)$ at generic $k$ by
Proposition~\ref{prop:pbw-universality} (PBW filtration).
$\mathrm{Vir}_c$ at generic $c$ by explicit bar computation.
$\mathcal{W}_N$ by DS descent from $V_k(\fg)$
(Proposition~\ref{prop:ds-koszul-descent}).
$V_\Lambda$ by the lattice Koszul theorem
(Theorem~\ref{thm:lattice-koszul}).
\end{proof}

\begin{proposition}[Explicit shadow free energies for the
standard families; \ClaimStatusProvedHere]%
\label{prop:g1-explicit-shadow-families}%
\index{shadow free energy!standard families}%
\begin{center}
\renewcommand{\arraystretch}{1.4}
\begin{tabular}{c|c|c}
$\cA$ & $\mathcal{F}^{\mathrm{sh}}(\cA;\hbar)$ &
  $Z^{\mathrm{sh}}(\cA;\hbar)$
\\[2pt]\hline\hline
$\mathcal{H}_k$ &
  $k\log(\hbar/2/\sinh(\hbar/2))$ &
  $(\hbar/2/\sinh(\hbar/2))^k$
\\
$V_k(\mathfrak{sl}_2)$ &
  $\frac{3(k+2)}{4}\log(\cdots)$ &
  $(\cdots)^{3(k+2)/4}$
\\
$\mathrm{Vir}_c$ &
  $\frac{5c}{6}\log(\cdots)
  + \sum_{r \geq 3} \mathcal{F}^{(r)}$ &
  $(\cdots)^{5c/6} \cdot e^{\sum \mathcal{F}^{(r)}}$
\\
$V_\Lambda$ &
  $\rank(\Lambda)\log(\cdots)$ &
  $(\cdots)^{\rank(\Lambda)}$
\end{tabular}
\end{center}
where $(\cdots) = \hbar/2/\sinh(\hbar/2)$.  The Virasoro
free energy receives nonlinear corrections at all arities
$r \geq 3$ because $r_{\max} = \infty$.
\end{proposition}

\begin{proof}
$\mathcal{F}^{\mathrm{sh}} = \kappa \cdot
\log(\hbar/2/\sinh(\hbar/2))$ by
Corollary~\ref{cor:g1-explicit-free-energy}.  The $\kappa$
values are: $k$ (Heisenberg), $(k+h^\vee)\dim\fg/(2h^\vee)$
(affine), $5c/6$ (Virasoro), $\rank(\Lambda)$ (lattice).
For Virasoro and $\mathcal{W}_N$, the nonlinear corrections
$\mathcal{F}^{(r)}$ are determined by the shadow tower
(Proposition~\ref{prop:g1-nonlinear-corrections}).
\end{proof}

\begin{proposition}[Virasoro: arity-$r$ amplitude bound;
\ClaimStatusProvedHere]%
\label{prop:g1-virasoro-bound}%
\index{Virasoro!arity bound}%
For $\mathrm{Vir}_c$ at generic $c$ with $|c| \geq c_0 > 0$
and $|5c + 22| \geq \delta_0 > 0$:
\begin{equation}\label{eq:g1-virasoro-bound}
|\mathcal{Z}_{r,g}(\mathrm{Vir}_c)|
\;\leq\;
\frac{K^r \cdot r!^3}{|c|^{r/2} |5c{+}22|^{r/4}}
\cdot q^{r/4},
\end{equation}
where $K$ is an absolute constant.  The arity series converges
for all $c \notin \{0, -22/5\}$.
\end{proposition}

\begin{proof}
The OPE coefficients of $\mathrm{Vir}_c$ are rational in $c$
with poles at $c = 0$ and $c = -22/5$.  The vertex norms satisfy
$\|m_k\|_q \leq K^k/(|c| \cdot |5c+22|)^{k/4}$ (the
Virasoro singular vectors grow as $(n+1)^4$ in the OPE).
The product over vertices at weight $r$ gives the denominator.
The HS-decay $q^{r/4}$ provides convergence.
\end{proof}

\begin{remark}[Admissible levels; \ClaimStatusConditional]%
\label{rem:g1-admissible}%
\index{admissible level!perturbative finiteness}%
At admissible levels $k + h^\vee = p/q$ with $p \geq h^\vee$,
the simple quotient $L_k(\fg)$ satisfies HS-sewing but may fail
Koszulness (vacuum null vectors violate diagonal Ext vanishing).
The graph amplitudes are still individually finite
(Proposition~\ref{prop:g1-log-integrability} uses only the FM
compactification), but the shadow tower interpretation requires
Koszulness.
\end{remark}

% ======================================================================
%
%  STAGE 6: IMPLICATIONS FOR 3D QUANTUM GRAVITY
%
% ======================================================================

\subsection{Implications for 3d quantum gravity}
\label{subsec:thqg-I-3d-gravity}
\index{3d quantum gravity!perturbative finiteness|textbf}
\index{BTZ black hole!partition function}
\index{twisted holography!3d implications}

The holographic modular Koszul datum
$\mathcal{H}(T) = (\cA, \cA^!, \mathcal{C}, r(z),
\Theta_\cA, \nabla^{\mathrm{hol}})$
(Definition~\ref{def:holographic-modular-koszul-datum})
specializes, for a 3d holomorphic-topological theory with
boundary chiral algebra $\cA$, to a gravitational
partition function.

\begin{definition}[Gravitational partition function]%
\label{def:g1-grav-partition}%
\index{gravitational partition function|textbf}%
The \emph{gravitational partition function} of an HT
theory $T$ with boundary algebra $\cA$ at coupling $\hbar$
is
\begin{equation}\label{eq:g1-grav-partition}
Z_{\mathrm{grav}}(T;\hbar)
\;:=\;
\sum_{g \geq 0} \hbar^{2g-2}
\int_{\overline{\mathcal{M}}_g}
\langle \Theta_\cA^{(g)} \rangle,
\end{equation}
where $\langle - \rangle$ is the trace on the shadow
algebra.  The genus-$0$ term is the classical action;
genus-$g$ for $g \geq 1$ is the $g$-loop correction.
\end{definition}

\begin{theorem}[Finiteness of the gravitational partition
function; \ClaimStatusProvedHere]%
\label{thm:g1-grav-finiteness}%
\index{gravitational partition function!finiteness|textbf}%
For any HT theory $T$ whose boundary algebra $\cA$ is
Koszul and satisfies HS-sewing, the gravitational partition
function $Z_{\mathrm{grav}}(T;\hbar)$ is a well-defined
formal power series in $\hbar^2$, convergent for $|\hbar|$
sufficiently small.  Each coefficient is a finite integral
over $\overline{\mathcal{M}}_g$.  No UV renormalization is
required.
\end{theorem}

\begin{proof}
Theorem~\ref{thm:g1-perturbative-finiteness}(d):
the genus series converges.  The integration over
$\overline{\mathcal{M}}_g$ is a finite integral of a
smooth form (the shadow $\langle \Theta^{(g)} \rangle$ is
a smooth cohomology class on $\overline{\mathcal{M}}_g$).
\end{proof}

\begin{proposition}[BTZ partition function from the
Heisenberg datum; \ClaimStatusProvedHere]%
\label{prop:g1-btz}%
\index{BTZ black hole!from Heisenberg}%
\index{Heisenberg!BTZ partition function}%
For $T$ = free chiral boson with boundary algebra
$\mathcal{H}_k$, the gravitational partition function
at genus $1$ is the BTZ black hole partition function:
\begin{equation}\label{eq:g1-btz}
Z_{\mathrm{grav}}(\mathcal{H}_k;\tau)
\;=\;
\sum_{\gamma \in
\operatorname{SL}_2(\mathbb{Z})/\Gamma_\infty}
Z_1(\mathcal{H}_k;\gamma\tau)
\;=\;
\sum_{\gamma}
(\operatorname{Im}(\gamma\tau))^{-k/2}
|\eta(\gamma\tau)|^{-2k},
\end{equation}
where the sum is over cosets
$\operatorname{SL}_2(\mathbb{Z})/\Gamma_\infty$ with
$\Gamma_\infty = \{\pm \bigl(\begin{smallmatrix}
1 & n \\ 0 & 1 \end{smallmatrix}\bigr)\}$, corresponding to
the sum over classical BTZ geometries.
\end{proposition}

\begin{proof}
The genus-$1$ contribution is
$Z_1(\mathcal{H}_k;\tau)$ from
Corollary~\ref{cor:g1-heisenberg-genus1}.
The sum over $\operatorname{SL}_2(\mathbb{Z})$-images is the
gravitational path integral over all genus-$1$ handlebodies
with the given boundary torus $\partial M = T^2_\tau$.
Each coset $\gamma \in
\operatorname{SL}_2(\mathbb{Z})/\Gamma_\infty$ corresponds
to a distinct solid torus filling: the contractible cycle is
$\gamma \cdot (0,1)^T$.  The thermal $\operatorname{AdS}_3$
is $\gamma = \bigl(\begin{smallmatrix} 1 & 0 \\
0 & 1 \end{smallmatrix}\bigr)$; the BTZ black hole is
$\gamma = \bigl(\begin{smallmatrix} 0 & -1 \\
1 & 0 \end{smallmatrix}\bigr)$.

Convergence of the sum: for
$\gamma = \bigl(\begin{smallmatrix} a & b \\
c & d \end{smallmatrix}\bigr)$ with $c \neq 0$,
$\operatorname{Im}(\gamma\tau) =
\operatorname{Im}\tau/|c\tau + d|^2$, so
$|\eta(\gamma\tau)|^{-2k}$ decays as $|c|^{-k}$ for large
$|c|$.  The number of cosets with $|c| \leq C$ is $O(C^2)$,
so the sum converges for $k > 2$ (sufficient decay in the
cusp).
\end{proof}

\begin{proposition}[Modular invariance of the gravitational
partition function; \ClaimStatusProvedHere]%
\label{prop:g1-modular-invariance}%
\index{gravitational partition function!modular invariance}%
$Z_{\mathrm{grav}}(T;\tau)$ is invariant under
$\operatorname{SL}_2(\mathbb{Z})$ acting on the boundary
modulus $\tau$: the sum over fillings
restores the modular invariance that the individual genus-$1$
contribution $Z_1(\cA;\tau)$ breaks.
\end{proposition}

\begin{proof}
The sum over $\operatorname{SL}_2(\mathbb{Z})/\Gamma_\infty$
is invariant under left $\operatorname{SL}_2(\mathbb{Z})$
action by construction: replacing $\tau \mapsto \sigma\tau$
permutes the cosets $\gamma \mapsto \gamma\sigma^{-1}$.
\end{proof}

\begin{remark}[The Farey tail and non-perturbative effects]%
\label{rem:g1-farey-tail}%
\index{Farey tail}%
\index{non-perturbative!effects}%
The Poincar\'e sum~\eqref{eq:g1-btz} is the genus-$1$
gravitational partition function.  At higher genus, the sum
over handlebody fillings of a genus-$g$ boundary involves
the mapping class group $\operatorname{MCG}(\Sigma_g)$ action
on the Teichm\"uller space $\mathcal{T}_g$, and the analog of
\eqref{eq:g1-btz} is
$Z_{\mathrm{grav}}^{(g)}(\cA;\Omega) =
\sum_{\gamma \in \operatorname{MCG}(\Sigma_g)/\Gamma_g}
Z_g(\cA;\gamma \cdot \Omega)$,
where $\Gamma_g$ is the stabilizer of a handlebody filling.
This sum is absolutely convergent under HS-sewing
(Theorem~\ref{thm:g1-perturbative-finiteness}).

Non-perturbative effects---contributions from non-handlebody
3-manifolds with the same boundary---are not captured by the
perturbative genus expansion.  They correspond, in the MC
framework, to non-perturbative MC elements that are not
accessible from the shadow tower.
\end{remark}

\begin{theorem}[Summary: perturbative finiteness of the
holographic modular Koszul datum; \ClaimStatusProvedHere]%
\label{thm:g1-finiteness-summary}%
\index{holographic modular Koszul datum!finiteness|textbf}%
For the holographic modular Koszul datum
$\mathcal{H}(T) = (\cA, \cA^!, \mathcal{C}, r(z),
\Theta_\cA, \nabla^{\mathrm{hol}})$
of an HT theory $T$ with Koszul boundary algebra satisfying
HS-sewing:
\begin{enumerate}[label=\textup{(\alph*)}]
\item Every component is perturbatively finite.
  $r(z) = \Res^{\mathrm{coll}}_{0,2}(\Theta_\cA)$ is a
  finite residue extraction.
  $\Theta_\cA$ has finite amplitudes at each $(r,g)$.
  $\nabla^{\mathrm{hol}}$ has finite coefficients at each
  order.
\item The MC equation
  $D_\cA\Theta_\cA + \frac{1}{2}[\Theta_\cA,\Theta_\cA] = 0$
  provides exact cancellation relations among the amplitudes
  at each bigraded component $(r,g)$:
  \begin{equation}\label{eq:g1-mc-cancellation}
  D\,\mathcal{Z}_{r,g}
  \;+\;
  \sum_{\substack{r_1+r_2=r+2 \\ g_1+g_2=g}}
  \tfrac{1}{2}[\mathcal{Z}_{r_1,g_1},
  \mathcal{Z}_{r_2,g_2}]
  \;+\;
  \Delta\,\mathcal{Z}_{r+2,g-1}
  \;=\; 0.
  \end{equation}
\item The FM compactification provides the intrinsic regulator.
  No counterterms, subtractions, or renormalization group
  actions are needed.
\item The entire standard landscape (Heisenberg, affine,
  Virasoro, $\mathcal{W}$-algebras, lattice VOAs) is covered.
\item The gravitational partition function is convergent for
  $|\hbar|$ sufficiently small.
\end{enumerate}
This is the foundational analytic result for the remaining
nine results
\textup{(G2)--(G10)} of this chapter.
\end{theorem}

\begin{proof}
(a): Theorem~\ref{thm:g1-perturbative-finiteness}(a)-(c) for
$\Theta_\cA$; the collision residue
$r(z) = \Res_{0,2}(\Theta_\cA)$ is a finite linear-algebraic
extraction from the finite arity-$2$ amplitude.  The shadow
connection $\nabla^{\mathrm{hol}} = d -
\operatorname{Sh}(\Theta_\cA)$ inherits finiteness from
$\Theta_\cA$.
(b): project the MC equation
$D_\cA\Theta + \frac{1}{2}[\Theta,\Theta] = 0$ onto
$(r,g)$-bigraded components; the differential $D_\cA =
\dzero + \sum_h \hbar^h d^{(h)}$ decomposes into tree-level
$\dzero$, BV operator $\Delta$, and higher-genus terms.
(c): Theorem~\ref{thm:g1-perturbative-finiteness}(e).
(d): Theorem~\ref{thm:g1-finiteness-landscape}.
(e): Theorem~\ref{thm:g1-grav-finiteness}.
\end{proof}

\begin{center}
\renewcommand{\arraystretch}{1.3}
\begin{tabular}{c|l}
\emph{Result} & \emph{Dependence on perturbative finiteness (G1)}
\\[2pt]\hline\hline
\textbf{(G2)} Complexity &
  Shadow depth $r_{\max}$ = arity of graph-sum termination \\
\textbf{(G3)} Polarization &
  Complementarity uses finite amplitudes on both summands \\
\textbf{(G4)} S-duality &
  $\kappa \leftrightarrow -\kappa$ acts on finite invariants \\
\textbf{(G5)} Yangian &
  $r(z) = \Res_{0,2}(\Theta_\cA)$ is a finite residue \\
\textbf{(G6)} Soft theorems &
  Ward identities are finite relations \\
\textbf{(G7)} Bootstrap &
  Genus spectral sequence has finite pages \\
\textbf{(G8)} Reconstruction &
  Algorithmic: finite input $\to$ finite output \\
\textbf{(G9)} Critical string &
  $c = 26$ self-duality uses finite amplitudes \\
\textbf{(G10)} Partition functions &
  Fredholm determinants from HS operators
\end{tabular}
\end{center}
